\begin{abstract}
	We propose the radiometric camera, a model that lifts any central
	geometric camera model to a radiometric one.  The lifted model
	keeps the projection of the geometric model and adds a finite
	aperture, a free-form focus surface and a generic vignetting
	function.  We calibrate it by inverse rendering directly on the
	relative irradiance of the sensor, from the same images of
	standard calibration targets used in geometric calibration.
	Geometric calibration absorbs the bias of its feature detections.
	The proposed radiometric calibration uses no features, and our
	radiometric cameras have lower reprojection errors on held-out
	test sets.  The focus surface extends defocus modeling from the
	thin lens to any central camera, and we use it in a first
	multi-view radiometric reconstruction, which recovers the in-focus
	reflectance of a texture from defocused images.  The
	vignetting is estimated jointly and accurately reproduces a flat
	field measured through a diffuser.  Because every standard camera
	model lifts to a radiometric one, the radiometric camera naturally
	and directly extends generative scene models beyond multi-view
	geometry to multi-view radiometric ones, where phenomena such as
	lighting and attenuation can be accurately estimated.
\end{abstract}
